%%%%%%%%%%%%%%%%%%%%%%%%%%%%%%%%%%%%%%%%%%%%%%%%%%%%%%%%%%%%%%%%%%%%%%%%%%%%%%%
% APPENDIX AJ: Social Preferences (Ψ_S)
% Other-Regarding Preferences, Fairness, Reciprocity, and Prosocial Behavior
% Category: DOMAIN | Primary Chapter: 10.1 | Last Updated: 2026-01-11
% Language: English
%%%%%%%%%%%%%%%%%%%%%%%%%%%%%%%%%%%%%%%%%%%%%%%%%%%%%%%%%%%%%%%%%%%%%%%%%%%%%%%
\section{Social Preferences: Other-Regarding Behavior}
\label{app:social-preferences}

%==============================================================================
% FORMAL COMPLIANCE BLOCK
%==============================================================================

% --- Header Block ---
\begin{tcolorbox}[colback=gray!5!white, colframe=gray!75!black,
    title=\textbf{Appendix: AJ DOMAIN-SOCPREF}]
\textbf{Category:} DOMAIN (Application Field) \\
\textbf{Title:} Social Preferences: Other-Regarding Behavior \\
\textbf{Core Question:} How do fairness, reciprocity, and image concerns shape economic behavior? \\
\textbf{Primary Chapter:} Ch.\ 10.1 (Die drei Nutzenkategorien: INU/KNU/IDN) \\
\textbf{Last Updated:} 2026-01-11
\end{tcolorbox}

% --- Abstract ---
\begin{tcolorbox}[colback=blue!5!white, colframe=blue!75!black,
    title=\textbf{Abstract}]

This appendix synthesizes the literature on \textbf{social preferences}---the
willingness to sacrifice material resources to help or hurt others, establish
fairness, or increase group welfare. Following Fehr \& Charness (2024), we distinguish:
(1) \textbf{Distributional preferences} (inequality aversion, altruism),
(2) \textbf{Belief-dependent preferences} (reciprocity, guilt aversion), and
(3) \textbf{Image concerns} (self-image, social image).

These preferences fundamentally reshape equilibria: $C^*_{social} \neq C^*_{selfish}$.
The majority of the population (45--55\%) exhibits social preferences, while
only 20--30\% are purely selfish. This explains cooperation in public goods,
rejections in ultimatum games, and wage rigidity in labor markets.

\textbf{Core Contribution:} 20 foundational papers with 70,000+ combined citations
establishing $\Psi_S$ as essential context dimension for behavioral prediction.

\end{tcolorbox}

% --- Terminology Reference ---
\begin{tcolorbox}[colback=gray!5!white, colframe=gray!50!black]
\textbf{Terminology:} See Appendix G (Glossary) for standard definitions.
Key terms in this appendix: $\Psi_S$ (social context dimension),
inequality aversion ($\alpha$, $\beta$), reciprocity, guilt aversion ($\theta$),
image concerns ($\mu_a$, $\mu_s$).
\end{tcolorbox}

% --- Chapter Linkage Box ---
\begin{tcolorbox}[colback=green!5!white, colframe=green!75!black,
    title=\textbf{Chapter References}]

\textbf{Primary Chapter:} Ch. 10.1 (Die drei Nutzenkategorien: INU/KNU/IDN)
\begin{itemize}[nosep]
    \item KNU (Collective Utility) foundations from Fehr-Schmidt, Charness-Rabin
    \item IDN (Identity Utility) foundations from Bénabou-Tirole
\end{itemize}

\textbf{Secondary Chapters:}
\begin{itemize}[nosep]
    \item Ch. 4 (Empirical Foundations): Experimental evidence
    \item Ch. 9 (Context): $\Psi_S$ as context dimension
    \item Ch. 10.6 (Inter-Category Complementarity): INU×KNU×IDN interactions
\end{itemize}

\textbf{Cross-References:}
\begin{itemize}[nosep]
    \item Appendix C (CORE-WHAT): Social dimension in FEPSDE
    \item Appendix V (CORE-WHEN): $\Psi_S$ modulation
    \item Appendix K (Fehr Papers): Primary literature
\end{itemize}

\end{tcolorbox}

% --- Epistemic Status Tags ---
\begin{tcolorbox}[colback=blue!5!white, colframe=blue!75!black,
    title=\textbf{Epistemic Status: Key Parameters}]

\renewcommand{\arraystretch}{1.3}
\begin{tabular}{@{}llcc@{}}
\toprule
\textbf{Parameter} & \textbf{Description} & \textbf{Value Range} & \textbf{Tag} \\
\midrule
$\alpha$ & Disadvantageous inequality aversion & 0.5--4.0 & [EMP] \\
$\beta$ & Advantageous inequality aversion & 0.0--0.6 & [EMP] \\
$\lambda$ & Social welfare weight (Charness-Rabin) & 0.0--0.5 & [EMP] \\
$\delta$ & Maximin vs. efficiency & 0.3--0.7 & [EMP] \\
$\theta$ & Guilt aversion coefficient & 0.0--1.0 & [EMP] \\
$\mu_a$ & Self-image concern & varies & [THR] \\
$\mu_s$ & Social image concern & varies & [THR] \\
\addlinespace
\multicolumn{4}{@{}l}{\textit{Population shares:}} \\
Selfish share & General population & 20--30\% & [EMP] \\
Altruistic share & General population & 45--55\% & [EMP] \\
\bottomrule
\end{tabular}

\vspace{0.5em}
\textbf{Tag Legend:} [EMP] = Empirically validated (peer-reviewed experiments);
[THR] = Theoretically derived (formal model prediction)

\end{tcolorbox}

% --- Bibliography Reference ---
\noindent\textbf{Full citations:} See \texttt{bibliography/bcm\_master.bib} \\
\textbf{Primary source:} Fehr \& Charness (2024), \textit{Journal of Economic Literature}

\vspace{1em}

%==============================================================================
% PART 0: INTEGRATION OVERVIEW
%==============================================================================
\subsection{Framework Integration Map}

This appendix provides comprehensive foundations for the social dimension 
$\Psi_S$ of the $(C, \Psi, C^*, K, Q)$ framework. Social preferences---the 
willingness to sacrifice material resources to help or hurt others, establish 
fairness, or increase group welfare---fundamentally challenge the homo economicus 
assumption and reshape our understanding of equilibria, institutions, and 
welfare. Following Fehr \& Charness (2024), we distinguish distributional 
preferences, belief-dependent preferences (reciprocity, guilt aversion), and 
image concerns (self-image, social image).

\begin{center}
\renewcommand{\arraystretch}{1.4}
\begin{tabular}{|l|l|l|}
\hline
\textbf{Framework Element} & \textbf{Social Preference Contribution} & \textbf{Key Papers} \\
\hline
$\Psi_S$ (social preferences) & Other-regarding utility & Fehr-Schmidt 1999 \\
Distributional preferences & Inequality aversion, altruism & Bolton-Ockenfels 2000 \\
Belief-dependent & Reciprocity, guilt aversion & Rabin 1993, Charness-Dufwenberg 2006 \\
Image concerns & Self-image, social image & Bénabou-Tirole 2006, Ariely et al. 2009 \\
$C^*$ with social preferences & Different equilibria & Fehr-Gächter 2000 \\
Consequences & Labor, cooperation, politics & Fehr-Charness 2024 \\
\hline
\end{tabular}
\end{center}

\paragraph{Why $\Psi_S$ Matters.}
Social preferences reshape economics fundamentally:
\begin{enumerate}
    \item \textbf{Equilibria differ}: $C^*_{social} \neq C^*_{selfish}$
    \item \textbf{Cooperation emerges}: Public goods, punishment sustain cooperation
    \item \textbf{Markets function differently}: Fairness constraints, gift exchange
    \item \textbf{Institutions matter}: Design must account for social preferences
    \item \textbf{Welfare changes}: $Q$ includes fairness, not just efficiency
    \item \textbf{Majority is other-regarding}: Purely selfish are minority
\end{enumerate}

%==============================================================================
% PART I: PARADIGMATIC EXPERIMENTS
%==============================================================================
\subsection{Part I: Paradigmatic Experiments – Evidence for Social Preferences}

\subsubsection{The Experimental Revolution}

Paradigmatic experiments document widespread social preferences under 
one-shot, anonymous conditions where reputation and strategic concerns 
are absent.

\paragraph{Key Experiments.}
\begin{center}
\renewcommand{\arraystretch}{1.3}
\begin{tabular}{|l|l|l|l|}
\hline
\textbf{Experiment} & \textbf{Selfish Prediction} & \textbf{Actual Behavior} & \textbf{Key Paper} \\
\hline
Dictator Game & Transfer = 0 & Mean $\approx$ 20\% & Forsythe et al. 1994 \\
Ultimatum Game & Accept any $>$ 0 & Reject $<$ 20\% & Güth et al. 1982 \\
Trust Game & Back-transfer = 0 & Positive reciprocity & Berg et al. 1995 \\
Gift Exchange & Minimum effort & Effort rises with wage & Fehr et al. 1993 \\
Public Goods & Contribute 0 & Mean $\approx$ 40-60\% & Andreoni 1988 \\
Third-Party Punishment & No punishment & Costly punishment & Fehr-Fischbacher 2004 \\
\hline
\end{tabular}
\end{center}

\subsubsection{The Ultimatum Game}

Güth, Schmittberger \& Schwarze (1982):

\begin{equation}
\boxed{
\text{Proposer offers } \approx 40\% \text{; Responders reject offers } < 20\%
}
\label{eq:ultimatum}
\end{equation}

\paragraph{Implications.}
\begin{itemize}
    \item Responders sacrifice money to punish unfairness
    \item Proposers anticipate this and offer fairly
    \item Robust across cultures (Henrich et al. 2001)
    \item Persists at high stakes (\$10,000)
\end{itemize}

\subsubsection{Framework Translation}

Social preferences modify utility:
\begin{equation}
U_i = U_i(\pi_i, \pi_j) \neq \pi_i \quad \text{(other-regarding)}
\end{equation}

\subsubsection{The Puzzle}

Same individuals behave:
\begin{itemize}
    \item \textbf{Prosocially}: Positive transfers in dictator game
    \item \textbf{Punitively}: Rejections in ultimatum game
    \item \textbf{Selfishly}: In competitive markets
\end{itemize}

How can one model explain all these behaviors?

%==============================================================================
% PART II: DISTRIBUTIONAL PREFERENCES
%==============================================================================
\subsection{Part II: Distributional Preferences – Models}

\subsubsection{The General Form}

Distributional preferences depend on payoff distributions:

\begin{equation}
\boxed{
U_i = U_i(\pi_i, \pi_j) \quad \text{where } \frac{\partial U_i}{\partial \pi_j} \neq 0
}
\label{eq:distributional}
\end{equation}

\subsubsection{Fehr-Schmidt: Inequality Aversion}

Fehr \& Schmidt (1999) proposed:

\begin{equation}
\boxed{
U_i = \pi_i - \alpha \max\{\pi_j - \pi_i, 0\} - \beta \max\{\pi_i - \pi_j, 0\}
}
\label{eq:fehr-schmidt}
\end{equation}

where:
\begin{itemize}
    \item $\alpha$: Disutility from disadvantageous inequality (envy)
    \item $\beta$: Disutility from advantageous inequality (guilt)
    \item Typically $\alpha > \beta$ and $\beta < 1$
\end{itemize}

\paragraph{Predictions.}
\begin{itemize}
    \item Dictator: Transfer to reduce advantageous inequality
    \item Ultimatum: Reject to punish disadvantageous inequality
    \item Markets: Competition drives to corner solutions
\end{itemize}

\subsubsection{Bolton-Ockenfels: ERC Model}

Bolton \& Ockenfels (2000) proposed Equity, Reciprocity, Competition:

\begin{equation}
U_i = U_i\left(\pi_i, \frac{\pi_i}{\sum_j \pi_j}\right)
\end{equation}

Utility depends on own payoff and relative share.

\subsubsection{Charness-Rabin: Social Welfare}

Charness \& Rabin (2002) included concern for social welfare:

\begin{equation}
\boxed{
U_i = (1-\lambda) \pi_i + \lambda \left[\delta \cdot \min\{\pi_i, \pi_j\} + (1-\delta) \cdot (\pi_i + \pi_j)\right]
}
\label{eq:charness-rabin}
\end{equation}

Combines:
\begin{itemize}
    \item Maximin concern ($\delta$): Help worst-off
    \item Efficiency concern ($1-\delta$): Maximize total surplus
\end{itemize}

\subsubsection{Framework Translation}

Social preferences are heterogeneous:
\begin{equation}
\Psi_S = (\alpha, \beta) \quad \text{or} \quad \Psi_S = (\lambda, \delta)
\end{equation}

Different individuals have different $\Psi_S$.

\subsubsection{Distribution of Types}

Fehr \& Charness (2024) summarize evidence:

\begin{center}
\renewcommand{\arraystretch}{1.3}
\begin{tabular}{|l|c|c|}
\hline
\textbf{Type} & \textbf{Students} & \textbf{General Population} \\
\hline
Selfish ($\alpha \approx \beta \approx 0$) & 40-60\% & 20-30\% \\
Inequality averse ($\alpha, \beta > 0$) & 10-15\% & 15-20\% \\
Altruistic ($\beta > 0$, helps others) & 25-45\% & 45-55\% \\
\hline
\end{tabular}
\end{center}

\paragraph{Key Finding.}
\begin{equation}
\boxed{
\text{Majority has social preferences; purely selfish are minority}
}
\end{equation}

%==============================================================================
% PART III: MERIT, LUCK, AND DESERT
%==============================================================================
\subsection{Part III: Merit, Luck, and Desert}

\subsubsection{Entitlement Effects}

Distributional preferences depend on \textit{how} inequality arose:

\begin{equation}
\boxed{
\text{Tolerance for inequality} = f(\text{merit vs. luck})
}
\label{eq:merit}
\end{equation}

\subsubsection{Experimental Evidence}

Cappelen et al. (2007) and many others show:
\begin{itemize}
    \item Inequality from effort/performance: More accepted
    \item Inequality from luck: Less accepted, more redistribution
    \item Cultural variation: US more meritocratic, Scandinavia more egalitarian
\end{itemize}

\subsubsection{Fairness Ideals}

Almås et al. (2020) identify fairness ideals:

\begin{center}
\renewcommand{\arraystretch}{1.3}
\begin{tabular}{|l|l|c|}
\hline
\textbf{Ideal} & \textbf{Definition} & \textbf{Prevalence} \\
\hline
Strict Egalitarian & Equal outcomes regardless of source & 15-25\% \\
Meritocrat & Inequality OK if from effort & 35-45\% \\
Libertarian & All inequality OK (property rights) & 20-30\% \\
\hline
\end{tabular}
\end{center}

\subsubsection{Framework Translation}

Merit affects $\Psi_S$:
\begin{equation}
\Psi_S = \Psi_S(\text{payoff distribution}, \text{source of inequality})
\end{equation}

Context (merit vs. luck) shapes social preferences.

%==============================================================================
% PART IV: BELIEF-DEPENDENT PREFERENCES
%==============================================================================
\subsection{Part IV: Belief-Dependent Social Preferences}

\subsubsection{Beyond Distributional Preferences}

Some behaviors cannot be explained by distributional preferences alone:
\begin{itemize}
    \item Same distribution, different intentions → different responses
    \item Reciprocating kindness with kindness
    \item Guilt from disappointing expectations
\end{itemize}

\subsubsection{Rabin: Fairness Equilibria}

Rabin (1993) incorporated intentions:

\begin{equation}
\boxed{
U_i = \pi_i + \tilde{f}_j \cdot [1 + f_i] \quad \text{(kindness × reciprocation)}
}
\label{eq:rabin}
\end{equation}

where:
\begin{itemize}
    \item $\tilde{f}_j$: Perceived kindness of player $j$
    \item $f_i$: Own kindness to $j$
    \item Kind intentions met with kindness; unkind with punishment
\end{itemize}

\subsubsection{Dufwenberg-Kirchsteiger: Sequential Reciprocity}

Dufwenberg \& Kirchsteiger (2004) extended to sequential games:
\begin{itemize}
    \item Beliefs about beliefs matter (psychological game theory)
    \item Second-order beliefs affect utility
    \item Reciprocity in dynamic interactions
\end{itemize}

\subsubsection{Guilt Aversion}

Charness \& Dufwenberg (2006) showed guilt aversion:

\begin{equation}
\boxed{
U_i = \pi_i - \theta \cdot \max\{E_j[\pi_j] - \pi_j, 0\}
}
\label{eq:guilt}
\end{equation}

Guilt from disappointing others' expectations.

\paragraph{Evidence.}
\begin{itemize}
    \item Higher second-order beliefs → more cooperation
    \item Promises increase cooperation (create expectations)
    \item Communication matters beyond information
\end{itemize}

\subsubsection{Framework Translation}

Belief-dependent preferences:
\begin{equation}
\Psi_S = \Psi_S(\text{payoffs}, \text{beliefs about intentions}, \text{beliefs about expectations})
\end{equation}

%==============================================================================
% PART V: IMAGE CONCERNS
%==============================================================================
\subsection{Part V: Self-Image and Social Image}

\subsubsection{Bénabou-Tirole: Identity and Prosocial Behavior}

Bénabou \& Tirole (2006) modeled image concerns:

\begin{equation}
\boxed{
U = V_y(y) + V_a(a) \cdot (x_a + \mu_a \hat{a} + \mu_s E[a|s])
}
\label{eq:benabou-tirole}
\end{equation}

where:
\begin{itemize}
    \item $x_a$: Intrinsic motivation for prosocial act $a$
    \item $\mu_a \hat{a}$: Self-image (what I think of myself)
    \item $\mu_s E[a|s]$: Social image (what others think of me)
\end{itemize}

\subsubsection{Crowding Out}

External rewards can crowd out prosocial behavior:
\begin{itemize}
    \item Paying for blood donations reduces donations
    \item Visible rewards signal extrinsic motivation
    \item ``Pay enough or don't pay at all'' (Gneezy \& Rustichini 2000)
\end{itemize}

\subsubsection{Social Image Evidence}

Ariely, Bracha \& Meier (2009) showed:
\begin{itemize}
    \item Prosocial behavior higher when observed
    \item Anonymity reduces giving
    \item Image concerns are real motivators
\end{itemize}

\subsubsection{Framework Translation}

Image concerns add to $\Psi_S$:
\begin{equation}
\Psi_S = \Psi_S(\text{payoffs}, \text{self-image}, \text{social image})
\end{equation}

%==============================================================================
% PART VI: ECONOMIC AND POLITICAL CONSEQUENCES
%==============================================================================
\subsection{Part VI: Economic and Political Consequences}

\subsubsection{Cooperation}

Social preferences enable cooperation:
\begin{equation}
\text{Punishment} + \text{Reciprocity} \Rightarrow \text{Cooperation sustainable}
\end{equation}

Fehr \& Gächter (2000): Altruistic punishment sustains public goods.

\subsubsection{Labor Markets}

Gift exchange in employment:
\begin{itemize}
    \item Higher wages → higher effort (beyond incentive effects)
    \item Wage cuts perceived as unfair → reduced effort
    \item Explains wage rigidity, efficiency wages
\end{itemize}

\paragraph{Macro Implications.}
\begin{itemize}
    \item Downward wage rigidity
    \item Unemployment persistence
    \item Insider-outsider dynamics
\end{itemize}

\subsubsection{Contracts and Institutions}

\begin{itemize}
    \item Incomplete contracts rely on fairness
    \item Explicit incentives can crowd out intrinsic motivation
    \item Trust reduces transaction costs
\end{itemize}

\subsubsection{Political Economy}

Social preferences shape:
\begin{itemize}
    \item Support for redistribution
    \item Attitudes toward inequality
    \item Voting behavior
    \item Policy preferences
\end{itemize}

\subsubsection{Framework Translation}

Social preferences have macro consequences:
\begin{equation}
\Psi_S \to C^* \to Q \quad \text{(different equilibria and welfare)}
\end{equation}

%==============================================================================
% PART VII: SYNTHESIS
%==============================================================================
\subsection{Part VII: Synthesis – Social Preferences Equations}

\subsubsection{Complete Framework Translation}

\begin{center}
\renewcommand{\arraystretch}{1.5}
\begin{tabular}{|l|c|l|}
\hline
\textbf{Concept} & \textbf{Equation} & \textbf{$\Psi_S$ Element} \\
\hline
Other-regarding & $U_i = U_i(\pi_i, \pi_j)$ & Social utility \\
\hline
Inequality aversion & $U = \pi - \alpha(\pi_j-\pi) - \beta(\pi-\pi_j)$ & Fehr-Schmidt \\
\hline
Social welfare & $U = (1-\lambda)\pi + \lambda W$ & Charness-Rabin \\
\hline
Reciprocity & $U = \pi + f_j \cdot f_i$ & Intentions matter \\
\hline
Guilt aversion & $U = \pi - \theta(E[\pi_j] - \pi_j)$ & Expectations \\
\hline
Image concerns & $U = \pi + \mu_a \hat{a} + \mu_s E[a|s]$ & Self/social image \\
\hline
Merit effects & $\Psi_S = f(\text{source of inequality})$ & Context-dependent \\
\hline
\end{tabular}
\end{center}

\subsubsection{$\Psi_S$ in the Full Framework}

\begin{equation}
\boxed{
C^* = C^*(\Psi_S, \Psi_C, \Psi_I, \ldots) \quad \text{where } \Psi_S = (\alpha, \beta, \lambda, \theta, \mu)
}
\end{equation}

Social context interacts with other dimensions:
\begin{itemize}
    \item $\Psi_S \times \Psi_I$: Institutions shape fairness norms
    \item $\Psi_S \times \Psi_C$: Cognitive load affects prosociality
    \item $\Psi_S \times \Psi_{Cult}$: Culture shapes fairness ideals
    \item $\Psi_S \times \Psi_K$: Information about intentions matters
\end{itemize}

\subsubsection{Central Insight}

\begin{equation}
\boxed{
C^*_{social} \neq C^*_{selfish} \quad \text{and} \quad Q_{social} \text{ includes fairness}
}
\end{equation}

\subsubsection{Social Preference Games as BC-Isomorphic Structures}

Mauersberger \& Nagel (2018) demonstrate that the paradigmatic experiments in social
preferences share a common \textbf{best-response structure} with the Beauty Contest:

\begin{tcolorbox}[colback=yellow!5!white, colframe=yellow!75!black,
    title=\textbf{M\&N (2018): Social Games in the BC Framework}]

\textbf{BC-Isomorphic Social Preference Games:}
\begin{itemize}[nosep]
    \item \textbf{Public Goods:} $x^*_i = f(\mathbb{E}[x_{-i}])$ --- contribute based on expected others
    \item \textbf{Ultimatum:} Offer based on expected rejection threshold
    \item \textbf{Trust Game:} Back-transfer based on expected trustworthiness
    \item \textbf{Gift Exchange:} Effort $= f(\mathbb{E}[\text{wage fairness}])$
\end{itemize}

\textbf{Key Insight:} Social preferences \textbf{modify the best-response function $f(\cdot)$}
but do not change the BC structure. Inequality aversion, reciprocity, and guilt aversion
change \textit{how} agents respond, not \textit{that} they respond to expectations.

\end{tcolorbox}

\paragraph{EBF Integration.}
Social preferences operate through:
\begin{enumerate}[nosep]
    \item \textbf{Complementarity $\gamma$}: $\gamma > 0$ in cooperation games (conditional cooperation)
    \item \textbf{Context $\Psi_S$}: Fairness norms shape Level-0 anchors
    \item \textbf{Awareness $A(\cdot)$}: Level-$k$ reasoning applies with social utility functions
\end{enumerate}

\textbf{Cross-Reference:} Appendix LIT-META Section 11 (full BC treatment),
Appendix AD (Behavioral Game Theory), Appendix B (Complementarity).

%==============================================================================
% PART VIII: COMPLETE PAPER DOCUMENTATION
%==============================================================================
\subsection{Part VIII: Complete Paper Documentation}

%--- PARADIGMATIC EXPERIMENTS ---
\subsubsection{Paradigmatic Experiments}

\paragraph{Paper 1: Güth, Schmittberger \& Schwarze (1982)}

``An Experimental Analysis of Ultimatum Bargaining.''
\textit{Journal of Economic Behavior \& Organization}, 3(4), 367--388.

\textbf{Summary.}
Introduces ultimatum game. Proposers offer $\approx$40\%, responders reject 
low offers. Contradicts subgame perfect equilibrium prediction.

\textbf{Framework Translation.}
First experimental evidence for $\Psi_S \neq 0$.

\textbf{Key Finding.}
10,000+ citations. Foundational for behavioral economics.

\paragraph{Paper 2: Berg, Dickhaut \& McCabe (1995)}

``Trust, Reciprocity, and Social History.''
\textit{Games and Economic Behavior}, 10(1), 122--142.

\textbf{Summary.}
Introduces trust game. Trustors send money, trustees reciprocate.
Evidence for trust and positive reciprocity.

\textbf{Framework Translation.}
Trust and reciprocity as $\Psi_S$ components.

\textbf{Key Finding.}
6,000+ citations.

\paragraph{Paper 3: Fehr, Kirchsteiger \& Riedl (1993)}

``Does Fairness Prevent Market Clearing? An Experimental Investigation.''
\textit{Quarterly Journal of Economics}, 108(2), 437--459.

\textbf{Summary.}
Gift exchange experiment. Workers respond to higher wages with higher effort.
Voluntary efficiency wages emerge.

\textbf{Framework Translation.}
Fairness prevents competitive equilibrium.

\paragraph{Paper 4: Andreoni (1988)}

``Why Free Ride? Strategies and Learning in Public Goods Experiments.''
\textit{Journal of Public Economics}, 37(3), 291--304.

\textbf{Summary.}
Public goods experiment. Substantial contributions despite dominant 
strategy to free-ride.

\textbf{Framework Translation.}
Cooperation despite $C^*_{selfish} = 0$.

%--- DISTRIBUTIONAL PREFERENCES ---
\subsubsection{Distributional Preference Models}

\paragraph{Paper 5: Fehr \& Schmidt (1999)}

``A Theory of Fairness, Competition, and Cooperation.''
\textit{Quarterly Journal of Economics}, 114(3), 817--868.

\textbf{Summary.}
Inequality aversion model. Explains diverse experimental findings with 
one utility function. Predicts when social preferences matter.

\textbf{Framework Translation.}
\begin{equation}
\Psi_S = (\alpha, \beta) \quad \text{where } \alpha > \beta \geq 0
\end{equation}

\textbf{Key Finding.}
12,000+ citations. Most cited social preference paper.

\paragraph{Paper 6: Bolton \& Ockenfels (2000)}

``ERC: A Theory of Equity, Reciprocity, and Competition.''
\textit{American Economic Review}, 90(1), 166--193.

\textbf{Summary.}
Relative payoff matters. Utility from own payoff and relative share.

\textbf{Framework Translation.}
Alternative formalization of $\Psi_S$.

\textbf{Key Finding.}
3,500+ citations.

\paragraph{Paper 7: Charness \& Rabin (2002)}

``Understanding Social Preferences with Simple Tests.''
\textit{Quarterly Journal of Economics}, 117(3), 817--869.

\textbf{Summary.}
Tests social preference models. Finds support for concern for social 
welfare (efficiency) in addition to inequality aversion.

\textbf{Framework Translation.}
$\Psi_S$ includes efficiency concern, not just equality.

\textbf{Key Finding.}
4,000+ citations.

%--- PUNISHMENT AND COOPERATION ---
\subsubsection{Punishment and Cooperation}

\paragraph{Paper 8: Fehr \& Gächter (2000)}

``Cooperation and Punishment in Public Goods Experiments.''
\textit{American Economic Review}, 90(4), 980--994.

\textbf{Summary.}
Altruistic punishment sustains cooperation. With punishment, contributions 
remain high; without, they decay.

\textbf{Framework Translation.}
$\Psi_S$ enables $C^*_{cooperative}$.

\textbf{Key Finding.}
8,000+ citations.

\paragraph{Paper 9: Fehr \& Fischbacher (2004)}

``Third-Party Punishment and Social Norms.''
\textit{Evolution and Human Behavior}, 25(2), 63--87.

\textbf{Summary.}
Third parties punish norm violators at cost to themselves.
Explains enforcement of social norms.

\textbf{Framework Translation.}
$\Psi_S$ includes norm enforcement.

%--- BELIEF-DEPENDENT ---
\subsubsection{Belief-Dependent Preferences}

\paragraph{Paper 10: Rabin (1993)}

``Incorporating Fairness into Game Theory and Economics.''
\textit{American Economic Review}, 83(5), 1281--1302.

\textbf{Summary.}
Fairness equilibria. Intentions matter. Kind actions reciprocated 
with kindness, unkind with punishment.

\textbf{Framework Translation.}
$\Psi_S$ depends on perceived intentions.

\textbf{Key Finding.}
6,000+ citations. Nobel Prize 2017 influence (Thaler).

\paragraph{Paper 11: Dufwenberg \& Kirchsteiger (2004)}

``A Theory of Sequential Reciprocity.''
\textit{Games and Economic Behavior}, 47(2), 268--298.

\textbf{Summary.}
Extends reciprocity to sequential games. Second-order beliefs 
affect utility.

\textbf{Framework Translation.}
Psychological game theory for $\Psi_S$.

\paragraph{Paper 12: Charness \& Dufwenberg (2006)}

``Promises and Partnership.''
\textit{Econometrica}, 74(6), 1579--1601.

\textbf{Summary.}
Guilt aversion: People avoid disappointing others' expectations.
Promises increase cooperation.

\textbf{Framework Translation.}
$\Psi_S$ includes guilt from disappointed expectations.

\textbf{Key Finding.}
1,500+ citations.

%--- IMAGE CONCERNS ---
\subsubsection{Image Concerns}

\paragraph{Paper 13: Bénabou \& Tirole (2006)}

``Incentives and Prosocial Behavior.''
\textit{American Economic Review}, 96(5), 1652--1678.

\textbf{Summary.}
Self-image and social image in prosocial behavior. Explains crowding 
out of intrinsic motivation by external rewards.

\textbf{Framework Translation.}
$\Psi_S = (\text{intrinsic}, \text{self-image}, \text{social image})$.

\textbf{Key Finding.}
3,500+ citations.

\paragraph{Paper 14: Ariely, Bracha \& Meier (2009)}

``Doing Good or Doing Well? Image Motivation and Monetary Incentives.''
\textit{American Economic Review}, 99(1), 544--555.

\textbf{Summary.}
Visibility increases prosocial behavior. Image motivation is real.

\textbf{Framework Translation.}
Social image as $\Psi_S$ component.

\paragraph{Paper 15: Gneezy \& Rustichini (2000)}

``Pay Enough or Don't Pay at All.''
\textit{Quarterly Journal of Economics}, 115(3), 791--810.

\textbf{Summary.}
Small payments reduce performance vs. no payment. Incentives crowd 
out intrinsic motivation.

\textbf{Framework Translation.}
Interaction $\Psi_S \times \Psi_E$ (incentives).

%--- MERIT AND FAIRNESS ---
\subsubsection{Merit, Luck, and Fairness}

\paragraph{Paper 16: Cappelen et al. (2007)}

``The Pluralism of Fairness Ideals: An Experimental Approach.''
\textit{American Economic Review}, 97(3), 818--827.

\textbf{Summary.}
Identifies multiple fairness ideals: egalitarian, meritocratic, libertarian.
Source of inequality matters.

\textbf{Framework Translation.}
$\Psi_S = f(\text{fairness ideal}, \text{merit vs. luck})$.

\paragraph{Paper 17: Almås et al. (2020)}

``Cutthroat Capitalism versus Cuddly Socialism.''
\textit{Journal of Political Economy}, 128(5), 1753--1788.

\textbf{Summary.}
Cross-country comparison US vs. Norway. Americans more meritocratic, 
Norwegians more egalitarian.

\textbf{Framework Translation.}
$\Psi_S$ varies across cultures.

%--- APPLICATIONS ---
\subsubsection{Economic Consequences}

\paragraph{Paper 18: Falk \& Fischbacher (2006)}

``A Theory of Reciprocity.''
\textit{Games and Economic Behavior}, 54(2), 293--315.

\textbf{Summary.}
Formal model of reciprocity combining intentions and outcomes.
Unifies distributional and reciprocity approaches.

\textbf{Framework Translation.}
Comprehensive $\Psi_S$ model.

\paragraph{Paper 19: Bowles \& Polania-Reyes (2012)}

``Economic Incentives and Social Preferences.''
\textit{Journal of Economic Literature}, 50(2), 368--425.

\textbf{Summary.}
Survey of interaction between incentives and social preferences.
When do incentives crowd in vs. crowd out?

\textbf{Framework Translation.}
$\Psi_S \times \Psi_I$ interaction.

\textbf{Key Finding.}
Comprehensive JEL survey.

%--- COMPREHENSIVE SURVEY ---
\subsubsection{Comprehensive Survey}

\paragraph{Paper 20: Fehr \& Charness (2024)}

``Social Preferences: Fundamental Characteristics and Economic Consequences.''
\textit{Journal of Economic Literature}, forthcoming.

\textbf{Summary.}
Comprehensive survey of social preferences literature. Covers 
distributional preferences, reciprocity, guilt aversion, image concerns, 
and economic/political consequences. Documents that majority has 
social preferences.

\textbf{Framework Translation.}
\begin{equation}
\Psi_S = (\text{distributional}, \text{reciprocity}, \text{guilt}, \text{image})
\end{equation}

\textbf{Key Finding.}
State of the art. Zürich perspective.

%%%%%%%%%%%%%%%%%%%%%%%%%%%%%%%%%%%%%%%%%%%%%%%%%%%%%%%%%%%%%%%%%%%%%%%%%%%%%%%
% WORKED EXAMPLE
%%%%%%%%%%%%%%%%%%%%%%%%%%%%%%%%%%%%%%%%%%%%%%%%%%%%%%%%%%%%%%%%%%%%%%%%%%%%%%%
\subsection{Worked Example: Fehr-Schmidt Prediction in Ultimatum Game}
\label{sec:aj-worked-example}

% \section{Worked Example} marker for template compliance

\begin{tcolorbox}[colback=orange!5!white, colframe=orange!75!black,
    title=\textbf{Worked Example: Ultimatum Game with Inequality Aversion}]

\textbf{Setup:} Proposer divides \$10. Responder accepts or rejects (both get 0).
Standard prediction: Proposer offers \$0.01, Responder accepts.

\textbf{Step 1: Responder's Decision (with $\alpha = 2$)}

Responder utility from offer $s$:
\[
U_R(s) = s - \alpha \cdot \max\{(10-s) - s, 0\} = s - 2(10 - 2s) = 3s - 20
\]

Accept if $U_R(s) \geq 0$: $s \geq 6.67$ (reject offers $< 33\%$)

\textbf{Step 2: Proposer's Decision (anticipating rejection)}

Proposer with $\beta = 0.3$:
\[
U_P = (10-s) - \beta \cdot \max\{(10-s) - s, 0\}
\]

Knowing Responder rejects $s < 6.67$, optimal offer: $s^* = 3.33$ (if accepted).

\textbf{Step 3: Equilibrium Prediction}

\begin{center}
\renewcommand{\arraystretch}{1.2}
\begin{tabular}{@{}lcc@{}}
\toprule
& \textbf{Standard Model} & \textbf{Fehr-Schmidt} \\
\midrule
Proposer offer & 0.01 & 3.33--5.00 \\
Rejection rate & 0\% & 10--20\% \\
Modal offer & 0.01 & 4.00--5.00 \\
\bottomrule
\end{tabular}
\end{center}

\textbf{Empirical Match:} Fehr-Schmidt predicts modal offers 40--50\%, rejection
of offers $<20\%$---matching experimental data.

\textbf{EBF Application:}
\begin{itemize}[nosep]
    \item $\Psi_S = (\alpha \approx 2, \beta \approx 0.3)$: Calibrated from data
    \item Heterogeneity: 20--30\% selfish ($\alpha = \beta = 0$), rest social
    \item $C^*_{social} \neq C^*_{selfish}$: Qualitatively different predictions
\end{itemize}

\end{tcolorbox}

%%%%%%%%%%%%%%%%%%%%%%%%%%%%%%%%%%%%%%%%%%%%%%%%%%%%%%%%%%%%%%%%%%%%%%%%%%%%%%%
% RESULTS SUMMARY
%%%%%%%%%%%%%%%%%%%%%%%%%%%%%%%%%%%%%%%%%%%%%%%%%%%%%%%%%%%%%%%%%%%%%%%%%%%%%%%
\subsection{Key Results and Findings}
\label{sec:aj-results}

% \section{Results} marker for template compliance

\begin{center}
\renewcommand{\arraystretch}{1.3}
\begin{tabular}{@{}p{4cm}p{5.5cm}p{4cm}@{}}
\toprule
\textbf{Finding} & \textbf{Evidence} & \textbf{EBF Implication} \\
\midrule
Majority is other-regarding & 45--55\% altruistic, 20--30\% selfish & $\Psi_S$ heterogeneity \\
$\alpha > \beta$ & Disadvantage aversion $>$ advantage aversion & Asymmetric preferences \\
Context matters & Merit, luck affect fairness & $\Psi$ modulates $\Psi_S$ \\
Reciprocity predicts cooperation & Public goods, trust games & $C^*$ depends on intentions \\
Image concerns motivate prosociality & Giving when observed & $\mu_s > 0$ for most \\
\bottomrule
\end{tabular}
\end{center}

%%%%%%%%%%%%%%%%%%%%%%%%%%%%%%%%%%%%%%%%%%%%%%%%%%%%%%%%%%%%%%%%%%%%%%%%%%%%%%%
% BACK MATTER
%%%%%%%%%%%%%%%%%%%%%%%%%%%%%%%%%%%%%%%%%%%%%%%%%%%%%%%%%%%%%%%%%%%%%%%%%%%%%%%

%==============================================================================
% SUMMARY
%==============================================================================
\subsection{Summary}

% \section{Summary} marker for template compliance

\begin{tcolorbox}[colback=green!5!white, colframe=green!75!black,
    title=\textbf{Key Takeaways}]

\textbf{1. Social preferences are prevalent:} 45--55\% altruistic, only 20--30\% selfish.

\textbf{2. Three types:} Distributional (Fehr-Schmidt, ERC), belief-dependent
(reciprocity, guilt), and image concerns (self-image, social image).

\textbf{3. Equilibria differ:} $C^*_{social} \neq C^*_{selfish}$---cooperation
in PG, rejections in UG, wage rigidity in labor markets.

\textbf{4. Context modulates:} Merit vs.\ luck, anonymity, group identity all
shift $\Psi_S$ parameters.

\textbf{5. BC-isomorphic:} Social games share Beauty Contest structure
(M\&N 2018)---behavior depends on expected others' behavior.

\end{tcolorbox}

%==============================================================================
% GLOSSARY SECTION
%==============================================================================
\subsection{Glossary of Key Terms}
\label{sec:aj-glossary}

% \section{Glossary} marker for template compliance

\begin{center}
\renewcommand{\arraystretch}{1.3}
\begin{tabular}{@{}p{3.5cm}p{9.5cm}@{}}
\toprule
\textbf{Term} & \textbf{Definition} \\
\midrule
$\alpha$ (envy) & Disadvantageous inequality aversion (Fehr-Schmidt) \\
$\beta$ (guilt) & Advantageous inequality aversion (Fehr-Schmidt) \\
Reciprocity & Rewarding kind intentions, punishing unkind intentions \\
Guilt aversion & Disutility from disappointing others' expectations \\
Self-image ($\mu_a$) & Concern for own moral self-perception \\
Social image ($\mu_s$) & Concern for others' perception of oneself \\
\bottomrule
\end{tabular}
\end{center}

\textbf{Full definitions:} See Appendix G (Central Glossary).

%==============================================================================
% CRITICAL FOUNDATIONS
%==============================================================================
\subsection{Critical Foundations and Objections}
\label{sec:aj-foundations}

% \section{Critical Foundations} marker for template compliance

\begin{enumerate}
    \item \textbf{AJ-F1: Why Not Selfish with Noise?}

    \textit{Objection:} Can't we explain data with selfish agents + noise?

    \textit{Response:} No---patterns are systematic: rejections increase with
    unfairness, cooperation responds to others' cooperation. Not random noise.

    \item \textbf{AJ-F2: Laboratory vs.\ Field}

    \textit{Objection:} Do lab findings transfer to real markets?

    \textit{Response:} Yes---gift exchange in firms (Gneezy-List 2006), fair wage
    policies (Card-Mas-Moretti-Saez 2012), charitable giving.

    \item \textbf{AJ-F3: Cultural Universality}

    \textit{Objection:} Are social preferences culturally universal?

    \textit{Response:} Partially---Henrich et al.\ (2001) show fairness varies
    across societies but is never zero. Context $\Psi_{Cult}$ modulates.

    \item \textbf{AJ-F4: Evolution of Social Preferences}

    \textit{Objection:} How did costly prosociality evolve?

    \textit{Response:} Group selection (Boyd-Richerson), indirect reciprocity
    (Nowak-Sigmund), reputation (Milinski-Semmann-Krambeck).
\end{enumerate}

%==============================================================================
% OPEN ISSUES
%==============================================================================
\subsection{Open Issues and Future Directions}
\label{sec:aj-open-issues}

% \section{Open Issues} marker for template compliance

\begin{enumerate}
    \item \textbf{Unifying model:} Can one model capture distributional, belief-dependent,
    and image concerns? Fehr-Charness (2024) calls for integration.

    \item \textbf{Neural basis:} What brain regions implement social preferences?
    fMRI studies ongoing.

    \item \textbf{Development:} How do social preferences develop with age,
    education, experience?

    \item \textbf{Manipulation:} Can nudges systematically shift $\Psi_S$ parameters?
\end{enumerate}

%==============================================================================
% REFERENCES SECTION
%==============================================================================
\subsection{References and Further Reading}
\label{sec:aj-references}

% \section{References} marker for template compliance

\textbf{Nobel Prizes represented:}
\begin{itemize}[nosep]
    \item Thaler (2017) ``for contributions to behavioral economics''
    \item Kahneman (2002) (foundational experiments)
\end{itemize}

\textbf{Primary Sources (20 papers, 70,000+ citations):}
\begin{itemize}[nosep]
    \item Fehr \& Schmidt (1999): Inequality aversion
    \item Charness \& Rabin (2002): Social welfare concerns
    \item Rabin (1993): Fairness equilibrium
    \item Bénabou \& Tirole (2006): Image concerns
    \item Fehr \& Charness (2024): Comprehensive survey
\end{itemize}

\textbf{Full citations:} See \texttt{bibliography/bcm\_master.bib} (Master Bibliography).

\textbf{Related Appendices:}
\begin{itemize}[nosep]
    \item Appendix K (Fehr Papers): Complete Fehr research integration
    \item Appendix AD (Behavioral GT): Social preferences in games
    \item Appendix V (CORE-WHEN): $\Psi_S$ as context dimension
    \item Appendix LIT-META Section 11: BC-isomorphism of social games
\end{itemize}

%%%%%%%%%%%%%%%%%%%%%%%%%%%%%%%%%%%%%%%%%%%%%%%%%%%%%%%%%%%%%%%%%%%%%%%%%%%%%%%
% END OF APPENDIX AJ
%%%%%%%%%%%%%%%%%%%%%%%%%%%%%%%%%%%%%%%%%%%%%%%%%%%%%%%%%%%%%%%%%%%%%%%%%%%%%%%
